% GENERATED FILE. Edit canonical lessons, then run npm run generate:books.
% canonical-source-hash: fnv1a64:9ad772c564b3a593
% canonical-lessons: ML-C182-ilakkuka, ML-C182-porikkuka, ML-C182-tilappikkuka, ML-C182-varukkuka, ML-C182-olikkuka

\chapter{To stir, To fry, To boil, To roast, To pour}
\label{ch:ml-to-stir-to-fry-to-boil-to-roast-to-pour}
\hlchaptermodality{182}

\begin{chapteropening}
\textbf{By the end of this chapter:} \emph{I can say to stir, to fry, to boil, to roast and to pour.}

Complete the last lesson of chapter 182: I can say to stir, to fry, to boil, to roast and to pour.
\end{chapteropening}

\section[\texorpdfstring{\ml{ഇളക്കുക} (iḷakkuka)}{ഇളക്കുക (iḷakkuka)}]{\ml{ഇളക്കുക} (iḷakkuka) — to stir}

\label{lesson:ML-C182-ilakkuka}



\begin{quote}
Before the new one: say the Malayalam for to pull, then the Malayalam for to push.
\end{quote}

\subsection*{You'll want to know: \ml{ഇളക്കുക}}
\textbf{\ml{ഇളക്കുക}} — \emph{iḷakkuka} — \textquotedblleft{}to stir\textquotedblright{}.

\begin{grammarlens}[title={What you've built}]
One more word for this chapter.
\end{grammarlens}

\subsection*{Your turn}
\begin{itemize}
\raggedright
  \item \emph{Say it:} \emph{iḷakkuka}
  \item \emph{Say it:} \emph{iḷakkuka}, once more
  \item \emph{Say it:} say \emph{iḷakkuka}
  \item \emph{Recall:} say the Malayalam for to pull, then the Malayalam for to push, then say \emph{iḷakkuka} again
\end{itemize}

\begin{tcolorbox}[breakable,colback=teal!4,colframe=teal!35!black,title={Before you move on}]
What does \ml{ഇളക്കുക} mean? (\textquotedblleft{}To stir\textquotedblright{}.) Say it once more.
\end{tcolorbox}
\section[\texorpdfstring{\ml{പൊരിക്കുക} (porikkuka)}{പൊരിക്കുക (porikkuka)}]{\ml{പൊരിക്കുക} (porikkuka) — to fry}

\label{lesson:ML-C182-porikkuka}



\begin{quote}
Before the new one: say the Malayalam for to push, then the Malayalam for to stir.
\end{quote}

\subsection*{You'll want to know: \ml{പൊരിക്കുക}}
\textbf{\ml{പൊരിക്കുക}} — \emph{porikkuka} — \textquotedblleft{}to fry\textquotedblright{}.

\begin{grammarlens}[title={What you've built}]
One more word for this chapter.
\end{grammarlens}

\subsection*{Your turn}
\begin{itemize}
\raggedright
  \item \emph{Say it:} \emph{porikkuka}
  \item \emph{Say it:} \emph{porikkuka}, once more
  \item \emph{Say it:} say \emph{porikkuka}
  \item \emph{Recall:} say the Malayalam for to push, then the Malayalam for to stir, then say \emph{porikkuka} again
\end{itemize}

\begin{tcolorbox}[breakable,colback=teal!4,colframe=teal!35!black,title={Before you move on}]
What does \ml{പൊരിക്കുക} mean? (\textquotedblleft{}To fry\textquotedblright{}.) Say it once more.
\end{tcolorbox}
\section[\texorpdfstring{\ml{തിളപ്പിക്കുക} (tiḷappikkuka)}{തിളപ്പിക്കുക (tiḷappikkuka)}]{\ml{തിളപ്പിക്കുക} (tiḷappikkuka) — to boil}

\label{lesson:ML-C182-tilappikkuka}



\begin{quote}
Before the new one: say the Malayalam for to stir, then the Malayalam for to fry.
\end{quote}

\subsection*{You'll want to know: \ml{തിളപ്പിക്കുക}}
\textbf{\ml{തിളപ്പിക്കുക}} — \emph{tiḷappikkuka} — \textquotedblleft{}to boil\textquotedblright{}.

\begin{grammarlens}[title={What you've built}]
One more word for this chapter.
\end{grammarlens}

\subsection*{Your turn}
\begin{itemize}
\raggedright
  \item \emph{Say it:} \emph{tiḷappikkuka}
  \item \emph{Say it:} \emph{tiḷappikkuka}, once more
  \item \emph{Say it:} say \emph{tiḷappikkuka}
  \item \emph{Recall:} say the Malayalam for to stir, then the Malayalam for to fry, then say \emph{tiḷappikkuka} again
\end{itemize}

\begin{tcolorbox}[breakable,colback=teal!4,colframe=teal!35!black,title={Before you move on}]
What does \ml{തിളപ്പിക്കുക} mean? (\textquotedblleft{}To boil\textquotedblright{}.) Say it once more.
\end{tcolorbox}
\section[\texorpdfstring{\ml{വറുക്കുക} (vaṟukkuka)}{വറുക്കുക (vaṟukkuka)}]{\ml{വറുക്കുക} (vaṟukkuka) — to roast}

\label{lesson:ML-C182-varukkuka}



\begin{quote}
Before the new one: say the Malayalam for to fry, then the Malayalam for to boil.
\end{quote}

\subsection*{You'll want to know: \ml{വറുക്കുക}}
\textbf{\ml{വറുക്കുക}} — \emph{vaṟukkuka} — \textquotedblleft{}to roast\textquotedblright{}.

\begin{grammarlens}[title={What you've built}]
One more word for this chapter.
\end{grammarlens}

\subsection*{Your turn}
\begin{itemize}
\raggedright
  \item \emph{Say it:} \emph{vaṟukkuka}
  \item \emph{Say it:} \emph{vaṟukkuka}, once more
  \item \emph{Say it:} say \emph{vaṟukkuka}
  \item \emph{Recall:} say the Malayalam for to fry, then the Malayalam for to boil, then say \emph{vaṟukkuka} again
\end{itemize}

\begin{tcolorbox}[breakable,colback=teal!4,colframe=teal!35!black,title={Before you move on}]
What does \ml{വറുക്കുക} mean? (\textquotedblleft{}To roast\textquotedblright{}.) Say it once more.
\end{tcolorbox}
\section[\texorpdfstring{\ml{ഒഴിക്കുക} (oḻikkuka)}{ഒഴിക്കുക (oḻikkuka)}]{\ml{ഒഴിക്കുക} (oḻikkuka) — to pour}

\label{lesson:ML-C182-olikkuka}



\begin{quote}
Before the new one: say the Malayalam for to boil, then the Malayalam for to roast.
\end{quote}

\subsection*{You'll want to know: \ml{ഒഴിക്കുക}}
\textbf{\ml{ഒഴിക്കുക}} — \emph{oḻikkuka} — \textquotedblleft{}to pour\textquotedblright{}.

\begin{grammarlens}[title={What you've built}]
One more word for this chapter.
\end{grammarlens}

\subsection*{Your turn}
\begin{itemize}
\raggedright
  \item \emph{Say it:} \emph{oḻikkuka}
  \item \emph{Say it:} \emph{oḻikkuka}, once more
  \item \emph{Say it:} say \emph{oḻikkuka}
  \item \emph{Recall:} say the Malayalam for to boil, then the Malayalam for to roast, then say \emph{oḻikkuka} again
\end{itemize}

\begin{tcolorbox}[breakable,colback=teal!4,colframe=teal!35!black,title={Before you move on}]
What does \ml{ഒഴിക്കുക} mean? (\textquotedblleft{}To pour\textquotedblright{}.) Say it once more.
\end{tcolorbox}
